\subsection*{S22. Retrospective source-surrogate architecture challenge}

This analysis asks whether the learner used to predict a fixed source target
changes source recovery and downstream hydration performance. It was designed
after the released-model and endpoint-architecture evaluations and is
exploratory, not a second prospective validation. The four challenged families
are SoluteML Abraham (8,098 source molecules), OpenFF (520), GBn2 (550) and
ConfSolv (17,829). CombiSolv-QM/COSMOtherm and MolSolv/SMD(water) retain their
released D-MPNN surrogates. Source target definitions, source rows, original
ARROW partitions, experimental endpoint labels and the 15-coordinate ordering
were fixed. OpenFF and GBn2 each have a calculated-energy target and an
experiment-minus-original-calculation residual target; the coordinate supplied
to the endpoint is the sum of the two predictions, not either target alone.

For each challenged family, a deterministic random 20\% source test was reserved
with seed 20261009. The remaining 80\% were divided into three shuffled
development folds with seed 20261010. Five candidates were specified before
source-test access: the released ExtraTrees or LightGBM architecture refitted
under the development splits; a residual MLP using the 217 RDKit descriptors
and binary Morgan bits; the same MLP using count-aware Morgan features;
count-aware ridge regression; and that ridge fit plus a bounded neural residual.
Neural fits used two width-128 residual blocks, dropout 0.1, AdamW at learning
rate 0.001, batch size 128, a 400-epoch ceiling and patience 35. A deterministic
10\% subset of each training fold selected the stopping epoch; the outer
source-development fold did not. Median imputation and feature scaling were
fitted on training rows only. Ridge used $\alpha=100$ on scaled RDKit and
log-count features. Neural targets were scaled within the training fold and
missing targets were masked in a Huber loss. Final source fits used the median
chosen development-fold epoch. The tree comparator preserved its released
target-specific leaf settings, and ConfSolv's variance-filtered bits were
chosen within each training fold.

Selection minimized the equal-target average of validation MAE divided by the
training standard deviation of each target, averaged over the three folds.
A precommitted difference smaller than 0.005 in this normalized score favored
the simpler candidate in the listed order. Thus selection used neither the
independent source test nor hydration test labels. An initial OpenFF development
screen used an incorrect tree comparator leaf size; it was invalidated before
source-test access, preserved in the experiment archive and rerun with the
released target-specific settings. The corrected screen is the one reported.
The selected models were the released tree for Abraham, ridge plus neural
residual for OpenFF and GBn2, and the released LightGBM tree for ConfSolv.
Supplementary Table~17 reports every independent source-test score. These
normalized errors compare candidate learners \emph{within} a family; different
source targets and units make a ranking across families inappropriate.

The selected hybrid's target-averaged source-test score was lower than the tree
score for OpenFF (0.2295 versus 0.2645) and GBn2 (0.2387 versus 0.2544).
Abraham's unselected hybrid happened to have the lowest source-test point
estimate (0.0948 versus 0.0973 for the selected tree), but the test did not
reverse the development-fold choice. ConfSolv's selected LightGBM tree also
had the lowest source-test point estimate (0.4930). A separate diagnostic held
out the largest 15\% by heavy-atom count within each source. The normalized
target-average score changed from 0.9002 to 0.7538 for OpenFF and from 1.0682
to 0.7118 for GBn2. Those scores again average separate targets. On the
independent random source test, the \emph{summed} corrected coordinate instead
changed from 0.660 to 0.651 kcal mol$^{-1}$ for OpenFF (selected-minus-tree
95\% paired interval, $-0.142$ to $0.128$) and from 0.718 to 0.683 for GBn2
($-0.204$ to $0.117$). On the size-separated source test, the corrected OpenFF
coordinate worsened from 1.896 to 2.282 kcal mol$^{-1}$ (interval, $-0.059$
to $0.847$), while GBn2 changed from 2.100 to 2.063 (interval, $-0.503$
to $0.441$). These intervals include zero. Component gains need not add when
errors in the two predicted terms are combined.

Source--endpoint identity overlap was audited by exact standardized molecule.
Every endpoint molecule appearing in a source table received source out-of-fold
predictions under that source condition. The cross-fit counts were 1,281 for
Abraham (1,189 public and 92 external), 520 for OpenFF (all public), 550 for
GBn2 (all public) and 151 for ConfSolv (139 public and 12 external). No ARROW-85
or Strict-97 molecule overlapped a supervised source table. This distinction is
especially important for the OpenFF and GBn2 experimental-residual targets:
they are not independent simulation-only supervision. External-220 has 123
source-exposed identities and is secondary to the strict source-disjoint 97
for the source-stage question.

Three response conditions were compared downstream: the original released
response matrix; \emph{source-tree-refit}, with all four challenged families
refitted using their released tree recipes and the same cross-fitting protocol;
and \emph{selected-all}, which substitutes only the validation-selected OpenFF
and GBn2 hybrids into that refit. The two unchanged D-MPNN families were checked
for exact equality across conditions. Each endpoint used identical labels,
outer folds, structural inputs, seed set and locked learner recipe between
source-tree-refit and selected-all. This produced 162 matched ExtraTrees fits
and 54 ridge-plus-neural-residual fits. Original-response flexible predictions
and released ExtraTrees results replayed their earlier values. Supplementary
Table~18 gives the six full condition--head results, rather than only favorable
substitutions. Paired differences use 10,000 molecule-level bootstrap resamples;
they are descriptive and not adjusted for the architecture search. For the
ExtraTrees head, selected-all minus source-tree-refit gives ARROW-85 $+0.0348$
(95\% interval, $-0.0009$ to $0.0771$), External-220 $-0.1175$
($-0.1765$ to $-0.0603$) and Strict-97 $-0.2428$
($-0.3585$ to $-0.1371$) kcal mol$^{-1}$. For the flexible head, the changes
are $+0.0038$, $-0.0041$ and $-0.0233$; all three intervals include zero.

The separate endpoint-size diagnostic holds out 192 public and 13 ARROW labels
under the earlier within-source size rule (Supplementary Method~S20). It is not
a strict global heavy-atom cutoff. Selected-all versus source-tree-refit changes
ExtraTrees MAE by $+0.1233$ kcal mol$^{-1}$ (interval, $0.0551$ to $0.1945$)
and the flexible endpoint by $-0.0141$ ($-0.0419$ to $0.0128$). Unlabelled
n-alkanes C2--C20 and Ac-(Gly/Ala)$_{0\text{--}6}$-NHMe were also scored. The
selected source hybrids continue changing on the alkane series beyond C10,
unlike the corresponding source trees, but the ExtraTrees endpoint remains
bounded and nearly plateaus on the peptide series. The flexible endpoint's
peptide predictions continue changing; no experimental peptide free energies
are available to establish which shape is accurate. Neither diagnostic proves
physical extensivity or conformational sensitivity. The released source
representation and public deployment artifact were not replaced as a result of
this retrospective comparison.
